Das Verfassen schriftlicher Ausarbeitungen ist faktisch in allen Studienfächern
unumgänglich. Neben einem fachlich fundierten Inhalt trägt auch eine
übersichtliche Gestaltung des Textbildes wesentlich zur Verständlichkeit einer
Arbeit bei. Diese Vorlage entstand, um die Hürden für die Benutzung des
Textsatzsystems \LaTeX\ zur Erstellung studentischer Arbeiten möglichst
niedrig zu legen. Aufgrund einer Vielzahl von unterschiedlichen Vorgaben, ist
diese Vorlage vielmehr eine als Bauskasten zu verstehende Sammlung von
Beispielen, als eine Layout-Vorgabe. Eigene Kommandos, werden nur an den
Stellen definiert, wo es unumgänglich ist (z.\,B. bei der Gestaltung der
Titelblätter etc.). Hierdurch ist es auch dem fortgeschrittenen
\LaTeX-Anwender möglich, eigene Anpassungen in diese Vorlage einzupflegen.

Kapitel~\ref{sec:vorlage} befasst sich mit dem Aufbau und der Funktionsweise
dieser Vorlage. Diese Kapitel sollte keinesfalls übersprungen werden, da es
dabei hilft, sich schnell im Grundgerüst dieser Vorlage zurecht zu finden.

\todo[inline]{Weitere Kapitel}

Wenn sich die Fertigstellung des Textkörpers abzeichnet, kann bereits die
formale Überarbeitung der Arbeit beginnen, indem ein Deckblatt und Inhalts- und
Abbildungsverzeichnisse hinzugefügt werden und bei fertigen Textpassagen
bereits die Rechtschreibkontrolle erfolgt. Gerade bei umfangereicheren Arbeiten
(Bachelorarbeit bzw. Diplomarbeit und aufwärts) erfolgt der Druck in einer
Druckerei. Für einen reibungslosen Ablauf sind einige Fallstricke zu beachten.
Kapitel~\ref{sec:final} fasst hierfür wichtig Hinweise zu diesem Thema
zusammen. Der Aufwand für die Endkontrolle kann durchaus zwei Wochen oder mehr
benötigen. Planen Sie also genügend zeitlichen Puffer ein und überspringen Sie
dieses Kapitel \textbf{nicht}, wenn es um's Ganze geht.


% Zielgruppe
% - Studenten
%   -> Projekt, Seminar-, Abschlussarbeiten
% - Promovenden, Habilitanten
%   -> Dissertationen, Habil.
%   -> Buchveröffentlichungen
